\documentclass[../../main.tex]{subfiles}

\begin{document}

\section{Conditionals and Loops}

Before getting started with loops, let's briefly go over increments.
For our counter $i$, we can \textbf{increment} or \textbf{decrement}
the value, or increase or decrease, by $1$ with \verb|i++| and
\verb|i--|.

For instance if we do the following,

\begin{basicbox}
\begin{lstlisting}[style=basic]
i = 5
i++
\end{lstlisting}
\end{basicbox}

then it is equivalent of the following statements.

\begin{basicbox}
\begin{lstlisting}[style=basic]
i = 5
i = i + 1
i = 6
\end{lstlisting}
\end{basicbox}

Similarly, \verb|i--| is identical to saying \verb|i = i - 1|.
There is also \verb|++i| and \verb|--i|, which first increments or
decrements the value before defining, but it is just better to
adjust index and use \verb|i++| and \verb|i--| to avoid confusion
unless there is a specific reason to use them.

\subsection{Conditionals}

One more thing to note before we move on is logical operations.
Unlike how we would usually write in math, we cannot use $=$ and
$\neq$ for equal to and not equal to. We also won't be able to use
comparisons $\geq$ and $\leq$. Instead, we would be using \verb|==|
for equal to, \verb|!=| for not equal to, \verb|<=| for less than or
equal to, and \verb|>=| for greater than or equal to. This is
crucial when we have to compare two values in different statements.
With these in mind, let's get started with if statements!

\subsubsection{If-Else Statements}

What if you want to display different statements depending on
the user input? Fortunately, we could achieve that through
if-else statements without manual human inspections. If, else,
if-else statements have the following forms.

\begin{basicbox}
\begin{lstlisting}[style=basic]
if (condition) {
    body
} else if (condition) {
    body
} else {
    body
}
\end{lstlisting}
\end{basicbox}

Let's take a look at a quick example.

\begin{cppbox}{CS0201.26051-01}{CS0201.26051-01}
\begin{lstlisting}[style=cppstyle]
#include <iostream>

int main() {
    int a;

    std::cout << "Enter a number: ";
    std::cin >> a;

    if (a > 10) {
        std::cout << "Your number is greater than ten!\n";
    } else if (a == 10) {
        std::cout << "Your number is ten!\n";
    } else {
        std::cout << "Your number is smaller than ten!\n";
    }

    return 0;
}
\end{lstlisting}
\end{cppbox}

In this short example, we compare an input from the user with $10$.
As you can see, we must use \verb|a == 10|, not \verb|a = 10|, as
it would be declaring $a$ as ten. We can see that our program works
as expected!

\begin{outputbox}
\begin{lstlisting}[style=basic]
$ ./program
Enter a number: -1
Your number is smaller than ten!

$ ./program
Enter a number: 10
Your number is ten!

$ ./program
Enter a number: 11
Your number is greater than ten!
\end{lstlisting}
\end{outputbox}

Instead of comparing numbers, we could also use statements for
comparing strings.

\begin{cppbox}{CS0201.26051-02}{CS0201.26051-02}
\begin{lstlisting}[style=cppstyle]
#include <iostream>

int main() {
    std::string password = "password";
    std::string input;

    std::cout << "Enter your password: ";
    std::cin >> input;

    bool correct = (input == password);

    if (correct) {
        std::cout << "Correct!\n";
    } else {
        std::cout << "Incorrect!\n";
    }

    return 0;
}
\end{lstlisting}
\end{cppbox}

As you can see from the example above we can take in the boolean
value. An alternative way would be directly putting
\verb|input == password| in the condition for if without defining
a new variable.

\begin{outputbox}
\begin{lstlisting}[style=basic]
$ ./program
Enter your password: password
Correct!

$ ./program
Enter your password: Password
Incorrect!
\end{lstlisting}
\end{outputbox}

If you have fixed options, and you don't want to type every
syntax for if-else statements, we can use \textbf{switch}
statements.

\subsubsection{Switch Statements}

Switch statements have the following structure.

\begin{basicbox}
\begin{lstlisting}[style=basic]
switch (expression) {
    case case1:
        block
    case case2:
        block
        break;
    default:
        block
}
\end{lstlisting}
\end{basicbox}

Here, \verb|break| and \verb|default| are optional, but nice to
have. Consider the following example for more intuitive explanation.

\begin{cppbox}{CS0201.26051-03}{CS0201.26051-03}
\begin{lstlisting}[style=cppstyle]
#include <iostream>

int main() {
    char input;

    std::cout << "Choices for Breakfast:\n"
              << "  A. Two Bananas\n"
              << "  B. Three Bananas\n"
              << "  C. Four Bananas\n"
              << "What would you like for breakfast today?\n";
    std::cin >> input;

    switch (input) {
        case 'A':
            std::cout << "You chose two bananas!\n";
            break;
        case 'B':
            std::cout << "You chose three bananas!\n";
            break;
        case 'C':
            std::cout << "You chose four bananas!\n";
            break;
        default:
            std::cout << "Please choose from the options\n";
    }

    return 0;
}
\end{lstlisting}
\end{cppbox}

Here, the program would go through all the options to see
if there are any matching cases. If not, it will execute
\verb|default| as the fallback. If you also want to exit from the
switch statement after finding the right case, we can do so
with the \verb|break| keyword.

\begin{outputbox}
\begin{lstlisting}[style=basic]
$ ./program
Choices for Breakfast:
  A. Two Bananas
  B. Three Bananas
  C. Four Bananas
What would you like for breakfast today?
B
You chose three bananas!

$ ./program
Choices for Breakfast:
  A. Two Bananas
  B. Three Bananas
  C. Four Bananas
What would you like for breakfast today?
D
Please choose from the options
\end{lstlisting}
\end{outputbox}

With the basic ideas of conditions in mind, let's discuss different
loops for C++.

\subsection{Loops}

There are many different fundamental loops that we can use in C++.
First, let's start with for loops.

\subsubsection{For Loops}

For loops retain the following structure.

\begin{basicbox}
\begin{lstlisting}[style=basic]
for (start; condition; update) {
    body
}
\end{lstlisting}
\end{basicbox}

The introduction of loops really made our lives easier. One of the
reason why we write program is to avoid repetitive tasks and work
more efficiently. Say we want to print ``I ate six bananas'' five
times on the screen. Without loops, we may have to use \verb|cout|
five times, which is ridiculous if the number increase to hundreds
and thousands. Instead, we can use a for loop.

\begin{cppbox}{CS0201.26051-04}{CS0201.26051-04}
\begin{lstlisting}[style=cppstyle]
#include <iostream>

int main() {
    for (int i = 0; i < 5; i++) {
        std::cout << "I ate six bananas\n";
    }

    return 0;
}
\end{lstlisting}
\end{cppbox}

When we compile and run, it is telling the computer to print
``I ate six bananas'' if $i < 5$ starting from $i = 0$ where it
increments by $1$ after each pass. As expected, we will get the
following result.

\begin{outputbox}
\begin{lstlisting}[style=basic]
I ate six bananas
I ate six bananas
I ate six bananas
I ate six bananas
I ate six bananas
\end{lstlisting}
\end{outputbox}

A more practical use of for loops would be this.

\begin{cppbox}{CS0201.26051-05}{CS0201.26051-05}
\begin{lstlisting}[style=cppstyle]
#include <iostream>

int main() {
    for (int i = 2; i <= 6; i++) {
        std::cout << "I ate " << i << " bananas\n";
    }

    return 0;
}
\end{lstlisting}
\end{cppbox}

Instead of manually changing the numbers, the loop will display
the following.

\begin{outputbox}
\begin{lstlisting}[style=basic]
I ate 2 bananas
I ate 3 bananas
I ate 4 bananas
I ate 5 bananas
I ate 6 bananas
\end{lstlisting}
\end{outputbox}

With this, it is natural to consider for loops inside a for loop.
This is called \textbf{nested loop} made with for loops. Consider
the following example that generates all possible combinations of
my breakfast assuming there exists three distinct bananas and
two distinct apple where I must eat exactly one banana and exactly
one apple.

\begin{cppbox}{CS0201.26051-06}{CS0201.26051-06}
\begin{lstlisting}[style=cppstyle]
#include <iostream>

int main() {
    for (int i = 0; i < 3; i++) {
        for (int j = 0; j < 2; j++) {
            std::cout << "Banana " << i+1 << " and "
                      << "Apple " << j+1 << std::endl;
        }
    }

    return 0;
}
\end{lstlisting}
\end{cppbox}

This will run the inner loop when $i = 0$, $i = 1$, and $i = 3$.
We also have to start from $i = 1$ and $j = 1$ or add the
displayed number by $1$, or else we will get Banana $0$ and
Apple $0$. Compiling and running, we can notice that we get
$3 \cdot 2 = 6$ distinct combinations.

\begin{outputbox}
\begin{lstlisting}[style=basic]
Banana 1 and Apple 1
Banana 1 and Apple 2
Banana 2 and Apple 1
Banana 2 and Apple 2
Banana 3 and Apple 1
Banana 3 and Apple 2
\end{lstlisting}
\end{outputbox}

Moving on from for loops, let's discuss while loops.

\subsubsection{While Loops}

While loops have the following form.

\begin{basicbox}
\begin{lstlisting}[style=basic]
while (condition) {
    body
}
\end{lstlisting}
\end{basicbox}

There are many things that we can do with while loops. Below are
two simple examples.

\begin{cppbox}{CS0201.26051-07}{CS0201.26051-07}
\begin{lstlisting}[style=cppstyle]
#include <iostream>

int main() {
    int i = 10;

    while (i > 0) {
        std::cout << i << " ";
        i--;
    }
    std::cout << std::endl;

    return 0;
}
\end{lstlisting}
\end{cppbox}

The code block inside the while loop will iterate until $i$ reaches
$1$.

\begin{outputbox}
\begin{lstlisting}[style=basic]
10 9 8 7 6 5 4 3 2 1
\end{lstlisting}
\end{outputbox}

Unlike the first example, the second example uses boolean for the
condition.

\begin{cppbox}{CS0201.26051-08}{CS0201.26051-08}
\begin{lstlisting}[style=cppstyle]
#include <iostream>

int main() {
    int birthdayYear = 9999;
    int input;
    bool wrong = true;

    while (wrong) {
        std::cout << "Guess which year I was born in!\n";
        std::cin >> input;

        if (input == birthdayYear) {
            wrong = false;
            std::cout << "Correct!\n";
        } else {
            std::cout << ":(\n";
        }
    }

    return 0;
}
\end{lstlisting}
\end{cppbox}

This program will ask the user the same question with feedback
until the boolean value \verb|wrong| is false, or the answer is
correct. Notice that it is important to put
\verb|std::cout << ":(\n";| in \verb|else| to prevent it showing
when the answer is correct.

\begin{outputbox}
\begin{lstlisting}[style=basic]
Guess which year I was born in!
1234
:(
Guess which year I was born in!
5678
:(
Guess which year I was born in!
9999
Correct!
\end{lstlisting}
\end{outputbox}

Lastly, below is an introduction to another loop.

\subsubsection{Do/While Loops}

Do/While loop is a variation of while loop that always run the first
time \cite{CS02-4}. Unlike the traditional while loop that always
check the condition first before running, do/while loop run the first
time, then check if next run is eligible even if the first pass does
not meet the condition. Below is the general structure of do/while
loops.

\begin{basicbox}
\begin{lstlisting}[style=basic]
do {
    block
} while (condition);
\end{lstlisting}
\end{basicbox}

Consider the following example.

\begin{cppbox}{CS0201.26051-09}{CS0201.26051-09}
\begin{lstlisting}[style=cppstyle]
#include <iostream>

int main() {
    bool correct = false;

    do {
        std::cout << "Correct!\n";
    } while (correct);

    return 0;
}
\end{lstlisting}
\end{cppbox}

With the traditional while loop, no outputs are expected. However,
do/while loop will give a different result.

\begin{outputbox}
\begin{lstlisting}[style=basic]
Correct!
\end{lstlisting}
\end{outputbox}

Notice that even if the condition is false for the first pass,
the block still runs anyhow.

This is it for the basics of conditionals and loops! Before discussing
fundamental algorithms, let's discuss more essential topics on
functions and references.

\end{document}


